% TOP_PROBLEM: {"id": 30006735, "title": "Cohen-Lenstra class-group distributions for quadratic number fields", "category_id": 1, "category": {"id": 1, "name": "number_theory", "display_name": "Number Theory", "description": "Properties of integers, prime numbers, Diophantine equations.", "slug": "number-theory", "order_index": 1, "created_at": "2026-07-31T15:26:25.670Z"}, "status": "open"}
% FIRST_POSTED: 2026-09-27
% !TeX program = lualatex
% Compile twice: lualatex 30006735.tex
% All references and prose are embedded; no BibTeX database or external inputs.
\documentclass[11pt,a4paper]{article}
\usepackage[margin=24mm,headheight=14pt]{geometry}
\usepackage{fontspec}
\setmainfont{Latin Modern Roman}
\setsansfont{Latin Modern Sans}
\setmonofont{Latin Modern Mono}
\usepackage{amsmath,amssymb,mathtools}
\usepackage{unicode-math}
\setmathfont{Latin Modern Math}
\usepackage{microtype}
\usepackage{enumitem,xurl,needspace}
\usepackage[dvipsnames]{xcolor}
\usepackage{hyperref}
\hypersetup{unicode=true,colorlinks=true,linkcolor=MidnightBlue,urlcolor=MidnightBlue,
 pdftitle={Research attempt: Problem 30006735},pdfauthor={Research notebook},bookmarksnumbered=true}
\usepackage{bookmark}
\usepackage{tocloft}
\setlength{\cftsecnumwidth}{3em}
\cftsetpnumwidth{2.5em}
\cftsetrmarg{4em}
\usepackage{titlesec}
\titleformat{\section}{\Large\bfseries\raggedright}{\thesection}{0.7em}{}
\usepackage{fancyhdr}
\pagestyle{fancy}\fancyhf{}
\fancyhead[L]{\small\sffamily Open Problems Math | Research notebook}
\fancyhead[R]{\small Problem 30006735 | 27 September 2026}
\fancyfoot[C]{\thepage}
\setlength{\parindent}{0pt}
\setlength{\parskip}{3pt plus 1pt}
\setlength{\emergencystretch}{3em}
\setcounter{tocdepth}{1}
\setcounter{secnumdepth}{1}
\setlist[itemize]{leftmargin=3em,itemsep=5pt,topsep=4pt,parsep=0pt}
\allowdisplaybreaks
\newcommand{\N}{\mathbb{N}}
\newcommand{\Z}{\mathbb{Z}}
\newcommand{\Q}{\mathbb{Q}}
\newcommand{\R}{\mathbb{R}}
\newcommand{\C}{\mathbb{C}}
\newcommand{\F}{\mathbb{F}}
\newcommand{\A}{\mathbb{A}}
\newcommand{\PP}{\mathbb P}
\newcommand{\E}{\mathbb{E}}
\newcommand{\eps}{\varepsilon}
\newcommand{\dd}{\,\mathrm d}
\newcommand{\sref}[2]{\hyperlink{source:#1:#2}{[S#2]}}
\newcommand{\eref}[2]{\hyperlink{extra:#1:#2}{[E#2]}}
% Turn the next switch off for a shorter reading copy. The quotations preserve
% every exactTarget field, including qualifications omitted from a short summary.
\newif\ifshowcatalogue
\showcataloguetrue
\newcommand{\cataloguescope}[1]{\ifshowcatalogue
 \paragraph{Exact scope in the supplied catalogue.}
 \begin{quote}\small #1\end{quote}\fi}
\usepackage{etoolbox}
\AtBeginEnvironment{itemize}{\small}
\numberwithin{equation}{section}
% Standalone export. Original research content preserved.
\begin{document}
\setcounter{section}{0}
% ================================================================
% problemId: 30006735
% Canonical problem ID: 30006735
% exactTarget SHA-256: e10166ca6c6246f53ba93e8d21a7e4e013b74a1c23f91774a75403d252b9a783
\clearpage
\section*{\texorpdfstring{Cohen-Lenstra class-group distributions for quadratic number fields}{Cohen-Lenstra class-group distributions for quadratic number fields}}\label{30006735}
\subsection{Definitions and mathematical statement}
The ideal class group $\operatorname{Cl}(K)$ is the group of nonzero fractional ideals modulo principal ideals. For an odd prime $p$, its $p$-primary subgroup consists of elements of $p$-power order. Let $\mathcal K_u(X)$ contain quadratic fields of discriminant absolute value below $X$, imaginary for $u=0$ and real for $u=1$. For every finite abelian $p$-group $A$, the claim is
\[
 \lim_{X\to\infty}\frac{\#\{K\in\mathcal K_u(X):\operatorname{Cl}(K)[p^\infty]\simeq A\}}
 {\#\mathcal K_u(X)}=\frac{|A|^{-u}}{Z_u(p)|\operatorname{Aut}(A)|},
\]
where $Z_u(p)=\sum_{[B]}|B|^{-u}/|\operatorname{Aut}(B)|$. This is a single-prime distribution for ordinary class groups, not a statement about $p=2$, narrow class groups, or joint independence across primes.
\par\smallskip\textit{Formulation and concept references:} \sref{30006735}{1}.
\cataloguescope{For u in \{0,1\} let K\_u(X) be the set of isomorphism classes of quadratic number fields K with |disc(K)| < X that are imaginary (u = 0) or real (u = 1). For each fixed odd prime p and each finite abelian p-group A, prove that the limit as X tends to infinity of \#\{K in K\_u(X) : Cl(K)[p\textasciicircum{}infinity] is isomorphic to A\} / \#K\_u(X) exists and equals (1/Z\_u(p)) * |A|\textasciicircum{}\{-u\} / |Aut(A)|, where Z\_u(p) is the sum over isomorphism classes of finite abelian p-groups B of |B|\textasciicircum{}\{-u\}/|Aut(B)| (so Z\_0(p) = prod\_\{i>=1\}(1-p\textasciicircum{}\{-i\})\textasciicircum{}\{-1\} and Z\_1(p) = prod\_\{i>=2\}(1-p\textasciicircum{}\{-i\})\textasciicircum{}\{-1\}). Cl(K) is the ordinary ideal class group and Cl(K)[p\textasciicircum{}infinity] its Sylow p-subgroup. No assertion is made for p = 2, for narrow class groups, for a joint law over several primes, or for any non-quadratic family.}
\subsection{Short English statement}
Do odd-primary class groups of quadratic number fields follow the exact automorphism-weighted distributions predicted by Cohen and Lenstra?
% BEGIN RESEARCH ADDITION 116
\subsection{Research attempt}

\paragraph{Current frontier and source audit.}
The source's introduction gives the two different weights for imaginary
and real quadratic fields, and records the Davenport--Heilbronn moments
$\E\#\operatorname{Sur}(\operatorname{Cl}(K),C_3)\to1$ and $1/3$,
respectively. These are not the full distributions of the $3$-primary
groups. The normalization incorporated by reference is checked in
Majumder's stated theorem. Throughout this attempt, $p$ is one fixed odd
prime and class groups are ordinary class groups. \sref{30006735}{1}\eref{30006735}{1}

Landesman--Levy's 2024 manuscript, revised in October 2025, obtains
surjection moments over quadratic \emph{function fields} for sufficiently
large finite constant fields. The permitted size of the constant field
depends on the target group. They explicitly warn that this does not
establish all moments at one fixed constant field. Neither that result
nor recent results about $2$-primary groups can be substituted for the
quadratic number-field target here. \eref{30006735}{2}

\paragraph{Attack route.}
The natural statistic is a surjection count rather than the order of the
class group. Three routes are examined: recover the law from all these
moments; derive finite, one-sided certificates for the probability of a
trivial $p$-part; and compute the same moments in a Haar random-matrix
model. The model calculation explains the proposed normalization but
supplies no arithmetic equidistribution theorem. The essential missing
input is stated explicitly rather than described as a generic randomness
assumption.

\paragraph{Attempt: an exact sufficient arithmetic input.}
Let $G_X=\operatorname{Cl}(K)[p^\infty]$ for a uniformly chosen
$K\in\mathcal K_u(X)$, and write
\[
 S_A(G)=\#\operatorname{Sur}(G,A),\qquad
 M_A(X)=\E S_A(G_X).
\]
Consider the additional hypothesis
\begin{equation}\label{eq:30006735:moments}
 \text{for every finite abelian $p$-group $A$,}\qquad
 M_A(X)\longrightarrow |A|^{-u}.
\end{equation}
This is a \emph{sufficient} condition, not something obtained in this
notebook for the arithmetic family. Pointwise convergence of group
probabilities, without tail control, must not be used to assert convergence
of these unbounded statistics in the reverse direction.

Here is a precise way to pass from \eqref{eq:30006735:moments} to the target.
Wood's Theorem~3.1 determines distributions modulo a fixed integer from
surjection moments with $M_A\leq|\bigwedge^2 A|$; Lemma~3.2 computes
those moments for the automorphism-weighted law. Our moments satisfy this
hypothesis because $|A|^{-u}\leq1\leq|\bigwedge^2 A|$.
For a desired group $A$ of exponent $p^e$, use modulus $p^{e+1}$.
A finite abelian $p$-group $G$ satisfies
\[
 G/p^{e+1}G\simeq A\quad\Longleftrightarrow\quad G\simeq A:
\]
a cyclic summand of length greater than $e$ would survive in the quotient
with length at least $e+1$. The cited fixed-modulus theorem therefore
gives exactly the asserted probability for $A$, not merely a distribution
of its rank. For $A=0$, take $e=0$. Apply the theorem along arbitrary
sequences $X\to\infty$ to obtain the full real-parameter limit.
\eref{30006735}{3}

\paragraph{Attempt: independently control escape and unbounded tests.}
The following elementary estimates expose two possible limit gaps.
Let $r(G)=\dim_{\F_p}G/pG$. Then
\[
 S_{C_p}(G)=p^{r(G)}-1,
 \qquad
 \limsup_{X\to\infty}\Pr(r(G_X)\geq R)
 \leq\frac{p^{-u}}{p^R-1}.
\]
If $\exp G\geq p^m$, its largest cyclic summand supplies at least
$\varphi(p^m)=p^{m-1}(p-1)$ surjections onto $C_{p^m}$. Hence
\begin{equation}\label{eq:30006735:tail}
 \limsup_{X\to\infty}\Pr(\exp G_X\geq p^m)
 \leq\frac{p^{-mu}}{p^{m-1}(p-1)}.
\end{equation}
Bounds on rank and exponent leave only finitely many group types. These
estimates imply tightness under \eqref{eq:30006735:moments}; no group-valued
probability mass escapes to arbitrarily large types.

For each fixed $A$, a pair of surjections to $A$ is a homomorphism to
$A\times A$, with additional surjectivity conditions. Partitioning all
homomorphisms by their image subgroup gives the exact identity and bound
\begin{equation}\label{eq:30006735:UI}
 S_A(G)^2\leq\#\operatorname{Hom}(G,A\times A)
 =\sum_{B\leq A\times A}S_B(G).
\end{equation}
The sum is over actual subgroups, so repetitions of an isomorphism type
are intentionally included. There are finitely many terms. Hypothesis
\eqref{eq:30006735:moments} bounds the expected right side, yielding uniformly
bounded second moments for $S_A(G_X)$ for large $X$. Therefore these
statistics are uniformly integrable. This separately justifies passing
their first moments to any subsequential limiting law. Merely having a
bounded first moment would not have justified that passage.

\paragraph{Attempt: finite-moment certificates for a trivial class-group part.}
There is a more explicit inversion for the probability of $G_X=0$.
For $r\geq0$ and $k\geq0$ let $\genfrac{[}{]}{0pt}{}{r}{k}_p$ be the Gaussian binomial
coefficient, zero when $k>r$. Since
\[
 S_{(C_p)^k}(G)=|\operatorname{GL}_k(\F_p)|\genfrac{[}{]}{0pt}{}{r(G)}{k}_p,
\]
define the finite observable
\[
 L_K(X)=\sum_{k=0}^K
 \frac{(-1)^k p^{k(k-1)/2}}{|\operatorname{GL}_k(\F_p)|}
 M_{(C_p)^k}(X),
\]
with $|\operatorname{GL}_0|=1$ and the zeroth moment equal to one.
Then, for every $X$ and $K\geq0$,
\begin{equation}\label{eq:30006735:bonferroni}
 L_{2K+1}(X)\leq\Pr(G_X=0)\leq L_{2K}(X).
\end{equation}
This holds for any random finite abelian $p$-group having the displayed
finite moments, without assuming Cohen--Lenstra. In particular, all of
these moments are finite for each finite set $\mathcal K_u(X)$ of quadratic number fields.

\emph{Proof.} The finite $q$-binomial identity gives, for $0\leq K<r$,
\[
 \sum_{k=0}^K(-1)^kp^{k(k-1)/2}\genfrac{[}{]}{0pt}{}{r}{k}_p
 =(-1)^Kp^{K(K+1)/2}\genfrac{[}{]}{0pt}{}{r-1}{K}_p.
\]
It follows by induction from the Gaussian-binomial recursion. When
$K\geq r>0$, the full sum is
$\prod_{j=0}^{r-1}(1-p^j)=0$; for $r=0$ it is one.
Thus each even partial sum is at least $\mathbf1_{r=0}$ and each odd
partial sum is at most it. Take expectations. A finite $p$-group has
rank zero exactly when it is trivial. $\square$

Under \eqref{eq:30006735:moments}, the limiting partial sums tend, as their
length increases, to
\begin{equation}\label{eq:30006735:euler}
 \sum_{k\geq0}
 \frac{(-1)^k p^{-k(k+1)/2-ku}}{\prod_{j=1}^k(1-p^{-j})}
 =\prod_{j\geq u+1}(1-p^{-j})=Z_u(p)^{-1}.
\end{equation}
The series is absolutely convergent: its denominator is bounded below by
$\prod_{j\geq1}(1-p^{-j})>0$, while its numerator decays quadratically in
$k$ in the exponent. The identity follows by expanding the convergent
product $\prod_{j\geq0}(1-zp^{-j})$ and putting $z=p^{-u-1}$.
Equation~\eqref{eq:30006735:bonferroni}, first with fixed $K$ and then with
$K\to\infty$, proves the limiting trivial-group probability without an
unjustified exchange of the $X$-limit and an infinite sum.

This also yields a quantitative finite-data test: estimated moments through
$(C_p)^{2K+1}$, with rigorous error bars, give an interval for the true
trivial-group probability. With exact moments the interval width is the final nonnegative term;
uncertainty in estimated moments must be added to the endpoints. That observation is an exact combinatorial consequence, not a
new estimate for arithmetic moments.

\paragraph{Attempt: derive the model moments instead of assuming them.}
Let $M_n:\Z_p^{n+u}\to\Z_p^n$ have independent Haar-uniform entries,
and $H_n=\operatorname{coker}M_n$. It is finite almost surely: a fixed
$n\times n$ minor has nonzero determinant almost surely, since the zero
set of a nonzero polynomial has Haar measure zero. For a fixed finite
abelian $p$-group $A$ of rank $r\leq n$, a surjection
$\phi:\Z_p^n\to A$ kills every column with probability
$|A|^{-(n+u)}$. Reduction modulo $p$ and the elementary full-rank matrix
count give
\[
 \#\operatorname{Sur}(\Z_p^n,A)
 =|A|^n\prod_{i=0}^{r-1}(1-p^{i-n}).
\]
Consequently
\begin{equation}\label{eq:30006735:Haar}
 \E S_A(H_n)=|A|^{-u}\prod_{i=0}^{r-1}(1-p^{i-n})
 \longrightarrow |A|^{-u}.
\end{equation}
For $r>n$ both surjection count and expectation are zero. This calculation
explains exactly how the rectangular excess $u$ enters. It does not show
that ideal-class relations behave like independent Haar columns; no
such independence is proved for quadratic fields. Wood's random-matrix
theorems rigorously address matrices, not that arithmetic identification.
\eref{30006735}{3}

\paragraph{Stress test.}
The first moment alone is radically insufficient. For either $u=0$ or $1$,
put mass $p^{-u}/(p-1)$ on $C_p$ and the remaining mass on the trivial
group. This is a probability law for odd $p$, and its $C_p$ moment is
exactly $p^{-u}$. But its $(C_p)^2$ moment is zero instead of $p^{-2u}$,
and all higher-exponent cyclic surjection moments vanish. Thus even a
correct Davenport--Heilbronn-type first moment cannot determine the
entire law.

The proposed route still needs \eqref{eq:30006735:moments} for \emph{every}
finite target group, in the prescribed number-field ordering. Estimates
for finitely many $A$, bounds without the exact main term, or a theorem
with a second varying field parameter do not supply it. The script checks
the Gaussian-binomial inequalities exactly and illustrates convergent
normalization sums; it does not compute or certify the required arithmetic
limits. No prediction for $p=2$, narrow class groups, or independence
across primes is appended.

\paragraph{Outcome.}
\textbf{Outcome: Conditional reduction.}

The result is a rigorous moment-to-law reduction with explicit rank and
exponent tails, a second-moment uniform-integrability bound, and finite
alternating-moment certificates for the trivial-group probability.
The exact Haar-model computation is separately established. These do not
constitute new arithmetic moment estimates, and no novelty claim is made
for the underlying moment method. Proving the missing number-field
surjection averages for all target groups would complete this route.
% END RESEARCH ADDITION 116
\subsection{Sources}
\begin{itemize}\raggedright
\item[\textnormal{[S1]}]\hypertarget{source:30006735:1}{} B. Alberts, Cohen-Lenstra Moments for Some Nonabelian Groups, JTNB 32 (2020), 631-664; normalization identity in P. Majumder, arXiv:1407.3066v1.
\url{https://www.numdam.org/item/10.5802/jtnb.1137.pdf}.
{\small\textit{Catalogue formulation source.}}
% BEGIN NEW REFERENCES 116
\item[\textnormal{[E1]}]\hypertarget{extra:30006735:1}{} Pritam Majumder,
\emph{An elementary proof of a power series identity for the weighted sum
of all finite abelian p-groups}, arXiv:1407.3066v1, 11 July 2014.
The theorem in Section 1 and Lemma 3 (automorphism counts).
\url{https://arxiv.org/pdf/1407.3066v1}.
\item[\textnormal{[E2]}]\hypertarget{extra:30006735:2}{} Aaron Landesman and Ishan Levy,
\emph{The Cohen--Lenstra moments over function fields via the stable
homology of non-splitting Hurwitz spaces}, arXiv:2410.22210v3,
1 October 2025. Introduction and Theorem 1.2.1; note the dependence of
the constant-field size on the target moment.
\url{https://arxiv.org/html/2410.22210v3}.
\item[\textnormal{[E3]}]\hypertarget{extra:30006735:3}{} Melanie Matchett Wood,
\emph{Random integral matrices and the Cohen--Lenstra Heuristics},
arXiv:1504.04391; Theorem 3.1, Lemma 3.2 and Corollaries 3.3--3.4
in the consulted manuscript.
\url{https://arxiv.org/pdf/1504.04391}.
% END NEW REFERENCES 116
\end{itemize}

\end{document}
